
\documentclass[12pt,a4paper]{ctexart}

\usepackage[margin=2.3cm]{geometry}
\usepackage{amsmath}
\usepackage{array}
\usepackage{booktabs}
\usepackage{caption}
\usepackage{enumitem}
\usepackage{etoolbox}
\usepackage{fancyhdr}
\usepackage{float}
\usepackage{graphicx}
\usepackage{hyperref}
\usepackage{longtable}
\usepackage{tabularx}
\usepackage[table]{xcolor}
\usepackage{tikz}
\usetikzlibrary{arrows.meta,positioning,shapes.geometric,fit,calc}

\hypersetup{hidelinks}
\urlstyle{same}
\setlist{nosep}
\captionsetup{font=small}
\renewcommand{\arraystretch}{1.25}
\setlength{\parindent}{2em}
\setlength{\parskip}{0.3em}
\setlength{\tabcolsep}{4pt}
\emergencystretch=3em
\sloppy
\makeatletter
\g@addto@macro{\UrlBreaks}{\do\/\do\-\do\_\do\.\do\:}
\makeatother
\AtBeginEnvironment{longtable}{\small\setlength{\tabcolsep}{3pt}}
\AtBeginEnvironment{tabularx}{\small\setlength{\tabcolsep}{3pt}}
\AtBeginEnvironment{tabular}{\small\setlength{\tabcolsep}{3pt}}

\newcommand{\studentname}{许奕}
\newcommand{\studentid}{2351441}
\newcommand{\teamname}{许奕、王天一、马源胜、周由}
\newcommand{\coursename}{软件工程管理与经济}
\newcommand{\projectname}{基于大模型的建筑能源智能管理与运维优化系统}
\newcommand{\docversion}{V2.1}
\newcommand{\docdate}{2026年6月}
\newcommand{\code}[1]{{\footnotesize\nolinkurl{#1}}}
\newcolumntype{Y}{>{\raggedright\arraybackslash}X}
\newcolumntype{C}[1]{>{\centering\arraybackslash}p{#1}}
\newcolumntype{L}[1]{>{\raggedright\arraybackslash}p{#1}}

\tikzset{
  box/.style={rectangle, rounded corners, draw=cyan!55!black, fill=cyan!8, thick, minimum height=8mm, align=center, inner sep=5pt},
  process/.style={rectangle, rounded corners, draw=teal!55!black, fill=teal!8, thick, minimum height=8mm, align=center, inner sep=5pt},
  data/.style={cylinder, shape border rotate=90, aspect=0.25, draw=orange!70!black, fill=orange!10, thick, minimum height=12mm, align=center, inner sep=5pt},
  risk/.style={diamond, draw=red!60!black, fill=red!8, thick, aspect=2, align=center, inner sep=2pt},
  arrow/.style={-{Stealth[length=2.4mm]}, thick, draw=gray!70!black},
  note/.style={rectangle, draw=gray!60, fill=gray!6, rounded corners, align=left, inner sep=6pt}
}

\pagestyle{fancy}
\fancyhf{}
\fancyhead[L]{\coursename}
\fancyhead[R]{\reporttitle}
\fancyfoot[C]{\thepage}
\setlength{\headheight}{15pt}

\title{
    \vspace{-2em}
    \heiti \zihao{2} \reporttitle \\
    \vspace{1em}
    \zihao{4} 课程期末项目正式交付文档
}
\author{
    \songti \zihao{4}
    项目名称：\projectname \\
    项目组成员：\teamname \\
    负责人：\studentname \quad 学号：\studentid
}
\date{\docdate}

\newcommand{\reporttitle}{SEE 软件经济分析与评价}
\begin{document}

\maketitle
\thispagestyle{fancy}

\begin{abstract}
本文档为《\projectname》的软件经济分析与评价（SEE），覆盖开发成本、维护成本、部署成本、定价策略、资金来源、收益测算、ROI、敏感性和综合经济评价。由于本项目为课程项目，文档同时采用课程实际投入口径和小规模校园试点口径，说明系统在课程交付和后续推广中的经济合理性。
\end{abstract}

\tableofcontents
\newpage

\section{文档控制}
\begin{longtable}{p{4cm}p{10.5cm}}
\toprule
项目 & 内容 \\
\midrule
文档名称 & SEE 软件经济分析与评价 \\
文档版本 & \docversion \\
评价范围 & 软件开发、部署运行、维护、定价、资金、收益、财务指标和风险敏感性 \\
评价对象 & 建筑能源智能管理与运维优化系统最终演示版本 \\
估算口径 & 课程项目影子成本 + 小规模校园试点假设 \\
\bottomrule
\end{longtable}

\section{经济评价范围}
\subsection{纳入范围}
本次经济评价纳入以下内容：
\begin{enumerate}
    \item 建筑能耗数据管理、查询、统计和导出能力。
    \item 后端 REST API、MCP Server、外部模型配置和知识库问答能力。
    \item 前端 Web 工作台，包括总览、数据浏览、统计分析、工单中心、决策报告和智能问答。
    \item 样例数据构造、数据字典、知识库、测试脚本、部署脚本和最终交付文档。
    \item 项目管理、测试、整合、系统展示和最终材料整理工作。
\end{enumerate}

\subsection{不纳入范围}
以下费用不纳入本次估算：真实传感器采购、楼宇自控系统接口费、生产级账号权限系统、正式商业合同税费、销售费用、法务费用和长期运维人员值班费用。

\section{开发成本估算}
\subsection{WBS 成本分解}
\begin{longtable}{p{4cm}p{5.8cm}p{2cm}p{2cm}}
\caption{工作分解结构与工时估算}\\
\toprule
工作包 & 主要内容 & 工时 & 占比 \\
\midrule
\endfirsthead
\toprule
工作包 & 主要内容 & 工时 & 占比 \\
\midrule
\endhead
需求与计划 & 作业要求分析、项目范围、任务拆解、接口契约 & 30 & 12.5\% \\
数据与知识库 & 数据源调研、样例数据生成、数据字典、知识库素材 & 42 & 17.5\% \\
后端开发 & FastAPI、分析服务、导出、工单、问答、MCP & 55 & 22.9\% \\
前端开发 & Vue 工作台、图表、三维视图、工单和交互 & 52 & 21.7\% \\
测试与联调 & pytest、前端构建、浏览器验收、修复问题 & 28 & 11.7\% \\
文档与提交 & SRS、SDS、SEE、SEM、验收、用户手册、PDF & 33 & 13.7\% \\
\midrule
合计 & 完成最终课程项目交付 & 240 & 100\% \\
\bottomrule
\end{longtable}

\subsection{成本构成图}
\begin{figure}[H]
\centering
\begin{tikzpicture}[x=0.055cm,y=0.42cm, every node/.style={font=\scriptsize}]
\draw[fill=cyan!25, draw=cyan!60!black] (0,0) rectangle (30,1);
\draw[fill=teal!25, draw=teal!60!black] (30,0) rectangle (72,1);
\draw[fill=orange!25, draw=orange!70!black] (72,0) rectangle (127,1);
\draw[fill=red!20, draw=red!60!black] (127,0) rectangle (179,1);
\draw[fill=purple!20, draw=purple!60!black] (179,0) rectangle (207,1);
\draw[fill=gray!20, draw=gray!60!black] (207,0) rectangle (240,1);
\node at (15,1.45) {需求 30h};
\node at (51,1.45) {数据 42h};
\node at (99.5,1.45) {后端 55h};
\node at (153,1.45) {前端 52h};
\node at (193,1.45) {测试 28h};
\node at (223.5,1.45) {文档 33h};
\node[align=center] at (120,-0.8) {总工时 240h，其中开发实现约占 44.6\%，测试和文档约占 25.4\%};
\end{tikzpicture}
\caption{课程项目工时构成}
\end{figure}

\subsection{人员成本}
按课程项目影子成本估算，人工单价取 80 元/小时。该单价不是实际工资，而是用于把学生投入转化为软件经济分析口径。

\begin{table}[H]
\centering
\caption{人员投入成本}
\begin{tabularx}{\textwidth}{p{2.5cm}Yp{2cm}p{2.4cm}}
\toprule
成员 & 主要职责 & 工时 & 成本 \\
\midrule
许奕 & 项目规划、整合、测试、文档、最终封版 & 80 & 6,400 元 \\
王天一 & 后端接口、架构、分析服务、MCP、测试 & 55 & 4,400 元 \\
马源胜 & 数据集、知识库、数据源调研、AI 配置 & 50 & 4,000 元 \\
周由 & 前端页面、图表、三维风险、交互联调 & 55 & 4,400 元 \\
\midrule
合计 & 需求、设计、开发、测试、文档、演示 & 240 & 19,200 元 \\
\bottomrule
\end{tabularx}
\end{table}

\subsection{工具与环境成本}
开发工具主要采用开源或免费方案。Python、FastAPI、Pandas、Vue、Vite、ECharts、GitHub 免费仓库均不产生直接软件授权费。考虑电脑折旧、网络和少量外部模型调试额度，课程阶段工具环境成本估算为 1,200 元。因此课程阶段一次性开发成本为：
\[
19,200 + 1,200 = 20,400 \text{ 元}
\]

\section{运行维护成本}
\subsection{本地演示成本}
课程验收支持本地运行和轻量云服务器演示两种形态。项目本身为数据分析与业务闭环系统，CPU 配置即可运行，不需要 GPU。轻量云演示主要用于让评审人员通过公网访问系统；本地演示仍可作为无公网环境下的回退方案。

\subsection{校园试点运行成本}
\begin{longtable}{p{4cm}p{2.8cm}p{2.8cm}p{4.4cm}}
\caption{校园试点年运行成本}\\
\toprule
项目 & 月费用 & 年费用 & 说明 \\
\midrule
\endfirsthead
\toprule
项目 & 月费用 & 年费用 & 说明 \\
\midrule
\endhead
轻量云服务器 & 100 元 & 1,200 元 & 2 核 4G 可支撑小规模试点，课程演示不需要 GPU \\
存储与备份 & 30 元 & 360 元 & 保存数据、日志和工单 \\
外部大模型 API & 100--500 元 & 1,200--6,000 元 & 可配置关闭，本地知识库兜底 \\
域名与 HTTPS & 50 元 & 600 元 & 校园内网部署可省略 \\
运维人力 & 2,160 元 & 25,920 元 & 数据检查、规则维护、故障处理和文档维护 \\
\midrule
合计 & 2,440--2,840 元 & 29,280--34,080 元 & 不含真实传感器与接口费 \\
\bottomrule
\end{longtable}

\section{收益分析}
\subsection{直接节能收益}
假设校园 4 栋建筑年电耗为 300,000 kWh，电价按 0.8 元/kWh 计算。系统通过异常识别、COP 监控、设备状态提醒和工单闭环，保守节能率取 5\%：
\[
\text{年节电量}=300,000 \times 5\% = 15,000 \text{ kWh}
\]
\[
\text{年节能收益}=15,000 \times 0.8 = 12,000 \text{ 元}
\]

\subsection{人工节约}
系统减少人工查表、人工定位异常和重复沟通。假设每个工作日节约 1 小时，全年 250 个工作日，人工机会成本 80 元/小时：
\[
1 \times 250 \times 80 = 20,000 \text{ 元/年}
\]

\subsection{异常损失减少}
异常设备长时间运行可能带来额外电费、设备磨损和服务投诉。保守假设每年避免 5 次中等异常，每次减少损失 2,000 元：
\[
5 \times 2,000 = 10,000 \text{ 元/年}
\]

\subsection{管理收益}
管理收益难以完全货币化，但对学校后勤管理和课程项目评审有重要价值：数据口径统一、异常定位更快、处理过程可追踪、报告自动生成、AI 客户端可调用系统能力、后续可扩展到真实部署。

\section{财务评价}
\subsection{课程项目口径}
课程项目一次性开发成本约 20,400 元。项目不产生直接收入，但形成完整软件资产、工程文档、演示系统、团队协作经验和后续复用能力。从课程评价角度，投入与成果匹配。

\subsection{四栋建筑试点口径}
四栋建筑试点情况下，年直接收益为：
\[
12,000 + 20,000 + 10,000 = 42,000 \text{ 元}
\]
若一次性实施成本取 25,000 元，年运行维护成本取 30,000 元，则第一年净收益为：
\[
42,000 - 25,000 - 30,000 = -13,000 \text{ 元}
\]
第一年 ROI 为：
\[
\frac{-13,000}{25,000 + 30,000} = -23.6\%
\]
第二年起不再承担一次性实施成本，年度净收益为：
\[
42,000 - 30,000 = 12,000 \text{ 元}
\]
第二年起 ROI 为：
\[
\frac{12,000}{30,000}=40\%
\]

\subsection{二十栋建筑扩展口径}
若扩展到 20 栋建筑，平台维护成本不会随建筑数量线性增长，而节能收益和异常减少收益会显著提升。假设年总收益 120,000 元，年运行维护成本 45,000 元，一次性实施成本 60,000 元，则：
\[
\text{年净收益}=120,000-45,000=75,000 \text{ 元}
\]
\[
\text{投资回收期}=\frac{60,000}{75,000}=0.8 \text{ 年}
\]
该结果说明系统在小范围试点时主要体现管理价值，在多建筑场景下具备更明确的经济收益。

\subsection{方案对比}
\begin{longtable}{L{3cm}L{3.2cm}L{3.2cm}L{3.2cm}L{2cm}}
\caption{不同部署规模的经济性对比}\\
\toprule
方案 & 一次性投入 & 年运行成本 & 年直接收益 & 评价 \\
\midrule
\endfirsthead
\toprule
方案 & 一次性投入 & 年运行成本 & 年直接收益 & 评价 \\
\midrule
\endhead
课程演示 & 20,400 元影子成本 & 近似 0 元 & 不产生直接收入 & 适合课程验收和能力展示 \\
4 栋试点 & 25,000 元 & 约 30,000 元 & 约 42,000 元 & 第二年起具备正向收益 \\
20 栋推广 & 60,000 元 & 约 45,000 元 & 约 120,000 元 & 回收期约 0.8 年 \\
生产级产品 & 视接口和权限复杂度确定 & 需加入安全、审计、运维成本 & 取决于客户数和节能率 & 适合后续产品化，不属于本期范围 \\
\bottomrule
\end{longtable}

\section{定价策略}
\subsection{项目制私有化部署}
适合学校、园区或企业内部使用。建议基础部署费 50,000 至 120,000 元，定制接口费按真实系统对接复杂度另计，年维护费按部署费 15\% 至 20\% 计取。优点是一次性收入明确，缺点是交付定制化压力较大。

\subsection{订阅制 SaaS}
适合后续产品化和多客户标准化使用。基础版可定价 1,000 至 3,000 元/月，专业版可定价 3,000 至 8,000 元/月，企业版按建筑数量、数据量、接口数量和模型调用量报价。该模式现金流稳定，但需要补充多租户、安全、权限和运维能力。

\subsection{咨询加平台混合模式}
当前项目成熟度更适合“能源诊断咨询 + 平台部署 + 后续维护”的混合模式。前期以数据整理和节能诊断服务进入，中期提供平台部署和培训，后期收取维护费、模型调用费和扩展开发费。

\section{资金与融资}
\begin{table}[H]
\centering
\caption{资金来源分析}
\begin{tabularx}{\textwidth}{p{3cm}YY}
\toprule
阶段 & 资金来源 & 说明 \\
\midrule
课程阶段 & 自有设备、免费开源工具、少量模型测试额度 & 不需要外部融资 \\
校内试点 & 学院实验教学经费、后勤节能管理经费、创新项目经费 & 适合验证系统价值 \\
推广阶段 & 节能减排专项、企业合作、节能服务公司分成 & 以真实节能收益支撑扩展 \\
商业化阶段 & 种子资金、产品化投入、渠道合作 & 用于多租户、安全和运维体系建设 \\
\bottomrule
\end{tabularx}
\end{table}

\section{敏感性分析}
\begin{figure}[H]
\centering
\begin{tikzpicture}[node distance=9mm and 10mm]
\node[box] (benefit) {经济收益};
\node[risk, above left=of benefit] (buildings) {建筑数量};
\node[risk, above right=of benefit] (saving) {节能率};
\node[risk, below left=of benefit] (maint) {维护成本};
\node[risk, below right=of benefit] (llm) {模型费用};
\draw[arrow] (buildings) -- node[above, sloped]{正相关} (benefit);
\draw[arrow] (saving) -- node[above, sloped]{正相关} (benefit);
\draw[arrow] (maint) -- node[below, sloped]{负相关} (benefit);
\draw[arrow] (llm) -- node[below, sloped]{负相关} (benefit);
\end{tikzpicture}
\caption{经济性敏感因素}
\end{figure}

敏感性结论如下：建筑数量和节能率是收益端最关键因素；维护成本和外部模型费用是成本端关键因素。若仅管理少数建筑，系统更偏管理工具；若扩展到多楼宇并形成统一运维流程，经济性明显提升。外部模型必须保持可关闭和可替换，以避免 API 费用成为不可控成本。

\section{综合评价}
从课程项目角度，本项目以较低现金成本形成完整软件系统和交付文档，体现了良好的学习与展示价值。从校园试点角度，系统能够改善数据管理、异常定位和运维闭环，短期经济收益不一定立刻覆盖全部投入，但管理价值明确。从多建筑推广角度，系统具备正向 ROI 和较短回收期。综合判断，本项目成本可控、收益路径清晰、推广潜力合理，符合课程中软件经济分析与评价要求。

\clearpage
\section{补充经济分析}
\subsection{工作包成本明细}
本节将 240 小时课程项目投入进一步分解到更细工作包，说明成本估算与需求、设计、开发、测试、文档和系统展示工作直接相关。
\begin{longtable}{L{1.5cm}L{3.2cm}L{5.8cm}L{1.6cm}L{2.4cm}}
\caption{工作包成本分解明细}\\
\toprule
编号 & 工作包 & 主要活动 & 工时 & 成本 \\
\midrule
\endfirsthead
\toprule
编号 & 工作包 & 主要活动 & 工时 & 成本 \\
\midrule
\endhead
W01 & 作业要求分析 & 阅读教师 PPT、整理最终交付范围、识别 SRS、SDS、SEE、SEM 和源码提交要求。 & 8 & 640 元 \\
W02 & 需求边界确定 & 确定建筑能源管理、异常分析、工单闭环、AI 问答和 MCP 的范围。 & 10 & 800 元 \\
W03 & 数据源调研 & 调研真实建筑能耗数据来源，形成模拟依据和字段口径。 & 8 & 640 元 \\
W04 & 样例数据生成 & 扩充多建筑、多楼层、多时段数据至 2026-06-01。 & 12 & 960 元 \\
W05 & 数据字典 & 整理字段、单位、派生字段和分析口径。 & 6 & 480 元 \\
W06 & 后端骨架 & 建立 FastAPI 应用、路由、配置和基础服务结构。 & 10 & 800 元 \\
W07 & 数据接口 & 实现概览、元信息、建筑、记录查询和导出接口。 & 10 & 800 元 \\
W08 & 统计服务 & 实现时间汇总、建筑对比、COP 排名和原因统计。 & 12 & 960 元 \\
W09 & 异常服务 & 实现规则识别、异常解释、楼层和设备分析。 & 14 & 1120 元 \\
W10 & 工单服务 & 实现创建、更新、排序和 JSON 持久化。 & 8 & 640 元 \\
W11 & 问答服务 & 实现知识库检索、外部模型配置和回退逻辑。 & 8 & 640 元 \\
W12 & MCP Server & 暴露 Tools、Resources 和 Prompt，复用服务层。 & 8 & 640 元 \\
W13 & 前端工作台 & 构建导航、状态管理、数据加载和页面布局。 & 10 & 800 元 \\
W14 & 数据浏览页 & 实现筛选、表格、导出和空状态。 & 6 & 480 元 \\
W15 & 统计分析页 & 实现范围筛选、异常优先、健康楼层和全局板块。 & 10 & 800 元 \\
W16 & 三维总览 & 实现 3D 楼宇、楼层颜色、拖拽旋转和点击联动。 & 10 & 800 元 \\
W17 & 工单中心 & 实现创建工单、完成工单、排序和持久化交互。 & 8 & 640 元 \\
W18 & 决策报告 & 实现运营日报、建议和收益模拟展示。 & 6 & 480 元 \\
W19 & AI 助手 & 实现模型选择、问答提交、回退和结果展示。 & 6 & 480 元 \\
W20 & 后端测试 & 编写接口、服务、工单、MCP 和 LLM 封装测试。 & 14 & 1120 元 \\
W21 & 前端构建验证 & 运行生产构建并修复构建错误。 & 6 & 480 元 \\
W22 & 浏览器验收 & 按演示路径检查页面、联动和刷新。 & 8 & 640 元 \\
W23 & SRS 文档 & 撰写需求、用例、功能表和追踪矩阵。 & 8 & 640 元 \\
W24 & SDS 文档 & 撰写架构、模块、数据、算法、状态机和接口设计。 & 8 & 640 元 \\
W25 & SEE 文档 & 完成成本、收益、定价、融资和经济评价。 & 6 & 480 元 \\
W26 & SEM 文档 & 完成范围、过程、质量、风险、配置和沟通管理。 & 6 & 480 元 \\
W27 & 用户手册 & 整理安装、启动、页面使用和故障排查。 & 4 & 320 元 \\
W28 & 提交材料 & 整理最终目录、PDF、源码副本和安全检查。 & 6 & 480 元 \\
\bottomrule
\end{longtable}

\subsection{收益与敏感性分析}
课程项目没有真实商业收费，因此经济分析采用“节能收益、人工节约、风险避免和管理价值”四类收益口径，并通过敏感性分析说明结果变化。
\begin{longtable}{L{2.4cm}L{3.0cm}L{5.6cm}L{3.8cm}}
\caption{收益与敏感性分析}\\
\toprule
情景 & 关键假设 & 测算说明 & 评价 \\
\midrule
\endfirsthead
\toprule
情景 & 关键假设 & 测算说明 & 评价 \\
\midrule
\endhead
保守节能 & 节电 2\% & 系统仅帮助发现明显夜间负荷和设备状态异常，收益较低但更可信。 & 适合作为最低收益边界。 \\
基准节能 & 节电 5\% & 结合异常定位、COP 优化和工单闭环，形成持续改进。 & 作为课程主估算口径。 \\
积极节能 & 节电 8\% & 运维团队及时执行建议并扩展到更多楼栋。 & 适合推广情景。 \\
人工节约 & 每周减少 4 小时巡检分析 & 系统自动汇总异常和报告，减少人工查表。 & 管理收益明显。 \\
风险避免 & 减少设备异常延迟发现 & 早发现低 COP、夜间负荷和设备状态异常，降低维修和投诉风险。 & 难完全量化但价值高。 \\
模型成本 & 外部模型可选 & 无 Key 时使用本地知识库，接入模型时按调用量产生边际成本。 & 成本可控。 \\
部署成本 & 本地运行与轻量云演示 & 课程演示可使用轻量云公网入口，后续生产推广再估算域名、HTTPS、备份和运维费用。 & 当前成本可控。 \\
数据成本 & 样例数据 & 当前使用模拟数据，真实部署需接入设备或能源平台。 & 推广时需新增预算。 \\
维护成本 & 每月少量维护 & 主要维护数据口径、异常规则、依赖升级和模型配置。 & 可由小团队承担。 \\
规模扩展 & 从 4 栋到 20 栋 & 服务层复用，主要增加数据接入和性能优化工作。 & 扩展性较好。 \\
失败情景 & 系统不被使用 & 若运维人员不按工单闭环处理，节能收益下降。 & 需管理制度配合。 \\
最佳情景 & 纳入月度运维流程 & 日报和工单成为固定流程，异常闭环形成管理证据。 & 经济价值最高。 \\
\bottomrule
\end{longtable}

\subsection{经济评价要素}
该表用于把课程要求的成本、收益、定价、融资和评价内容逐项落到项目本身，避免经济分析与系统功能脱节。
\begin{longtable}{L{1.8cm}L{3.3cm}L{6.2cm}L{4.0cm}}
\caption{经济评价要素}\\
\toprule
编号 & 评价项 & 分析说明 & 文档体现 \\
\midrule
\endfirsthead
\toprule
编号 & 评价项 & 分析说明 & 文档体现 \\
\midrule
\endhead
EC-01 & 需求分析投入 & 围绕“需求分析投入”说明其与建筑能源管理系统的关系，并给出成本、收益或管理价值判断。 & SEE 正文或补充表格中体现。 \\
EC-02 & 数据调研投入 & 围绕“数据调研投入”说明其与建筑能源管理系统的关系，并给出成本、收益或管理价值判断。 & SEE 正文或补充表格中体现。 \\
EC-03 & 后端接口投入 & 围绕“后端接口投入”说明其与建筑能源管理系统的关系，并给出成本、收益或管理价值判断。 & SEE 正文或补充表格中体现。 \\
EC-04 & 分析算法投入 & 围绕“分析算法投入”说明其与建筑能源管理系统的关系，并给出成本、收益或管理价值判断。 & SEE 正文或补充表格中体现。 \\
EC-05 & 前端页面投入 & 围绕“前端页面投入”说明其与建筑能源管理系统的关系，并给出成本、收益或管理价值判断。 & SEE 正文或补充表格中体现。 \\
EC-06 & 三维交互投入 & 围绕“三维交互投入”说明其与建筑能源管理系统的关系，并给出成本、收益或管理价值判断。 & SEE 正文或补充表格中体现。 \\
EC-07 & 工单功能投入 & 围绕“工单功能投入”说明其与建筑能源管理系统的关系，并给出成本、收益或管理价值判断。 & SEE 正文或补充表格中体现。 \\
EC-08 & AI 问答投入 & 围绕“AI 问答投入”说明其与建筑能源管理系统的关系，并给出成本、收益或管理价值判断。 & SEE 正文或补充表格中体现。 \\
EC-09 & MCP 集成投入 & 围绕“MCP 集成投入”说明其与建筑能源管理系统的关系，并给出成本、收益或管理价值判断。 & SEE 正文或补充表格中体现。 \\
EC-10 & 测试投入 & 围绕“测试投入”说明其与建筑能源管理系统的关系，并给出成本、收益或管理价值判断。 & SEE 正文或补充表格中体现。 \\
EC-11 & 文档投入 & 围绕“文档投入”说明其与建筑能源管理系统的关系，并给出成本、收益或管理价值判断。 & SEE 正文或补充表格中体现。 \\
EC-12 & 系统展示投入 & 围绕“系统展示投入”说明其与建筑能源管理系统的关系，并给出成本、收益或管理价值判断。 & SEE 正文或补充表格中体现。 \\
EC-13 & 人员单价 & 围绕“人员单价”说明其与建筑能源管理系统的关系，并给出成本、收益或管理价值判断。 & SEE 正文或补充表格中体现。 \\
EC-14 & 工具成本 & 围绕“工具成本”说明其与建筑能源管理系统的关系，并给出成本、收益或管理价值判断。 & SEE 正文或补充表格中体现。 \\
EC-15 & 部署成本 & 围绕“部署成本”说明其与建筑能源管理系统的关系，并给出成本、收益或管理价值判断。 & SEE 正文或补充表格中体现。 \\
EC-16 & 模型调用成本 & 围绕“模型调用成本”说明其与建筑能源管理系统的关系，并给出成本、收益或管理价值判断。 & SEE 正文或补充表格中体现。 \\
EC-17 & 维护成本 & 围绕“维护成本”说明其与建筑能源管理系统的关系，并给出成本、收益或管理价值判断。 & SEE 正文或补充表格中体现。 \\
EC-18 & 数据接入成本 & 围绕“数据接入成本”说明其与建筑能源管理系统的关系，并给出成本、收益或管理价值判断。 & SEE 正文或补充表格中体现。 \\
EC-19 & 节能收益 & 围绕“节能收益”说明其与建筑能源管理系统的关系，并给出成本、收益或管理价值判断。 & SEE 正文或补充表格中体现。 \\
EC-20 & 人工节约 & 围绕“人工节约”说明其与建筑能源管理系统的关系，并给出成本、收益或管理价值判断。 & SEE 正文或补充表格中体现。 \\
EC-21 & 风险避免 & 围绕“风险避免”说明其与建筑能源管理系统的关系，并给出成本、收益或管理价值判断。 & SEE 正文或补充表格中体现。 \\
EC-22 & 管理价值 & 围绕“管理价值”说明其与建筑能源管理系统的关系，并给出成本、收益或管理价值判断。 & SEE 正文或补充表格中体现。 \\
EC-23 & 课程评分价值 & 围绕“课程评分价值”说明其与建筑能源管理系统的关系，并给出成本、收益或管理价值判断。 & SEE 正文或补充表格中体现。 \\
EC-24 & 推广试点价值 & 围绕“推广试点价值”说明其与建筑能源管理系统的关系，并给出成本、收益或管理价值判断。 & SEE 正文或补充表格中体现。 \\
EC-25 & 保守情景 & 围绕“保守情景”说明其与建筑能源管理系统的关系，并给出成本、收益或管理价值判断。 & SEE 正文或补充表格中体现。 \\
EC-26 & 基准情景 & 围绕“基准情景”说明其与建筑能源管理系统的关系，并给出成本、收益或管理价值判断。 & SEE 正文或补充表格中体现。 \\
EC-27 & 积极情景 & 围绕“积极情景”说明其与建筑能源管理系统的关系，并给出成本、收益或管理价值判断。 & SEE 正文或补充表格中体现。 \\
EC-28 & 成本敏感性 & 围绕“成本敏感性”说明其与建筑能源管理系统的关系，并给出成本、收益或管理价值判断。 & SEE 正文或补充表格中体现。 \\
EC-29 & 收益敏感性 & 围绕“收益敏感性”说明其与建筑能源管理系统的关系，并给出成本、收益或管理价值判断。 & SEE 正文或补充表格中体现。 \\
EC-30 & 盈亏平衡 & 围绕“盈亏平衡”说明其与建筑能源管理系统的关系，并给出成本、收益或管理价值判断。 & SEE 正文或补充表格中体现。 \\
EC-31 & 定价策略 & 围绕“定价策略”说明其与建筑能源管理系统的关系，并给出成本、收益或管理价值判断。 & SEE 正文或补充表格中体现。 \\
EC-32 & 融资假设 & 围绕“融资假设”说明其与建筑能源管理系统的关系，并给出成本、收益或管理价值判断。 & SEE 正文或补充表格中体现。 \\
EC-33 & 维护预算 & 围绕“维护预算”说明其与建筑能源管理系统的关系，并给出成本、收益或管理价值判断。 & SEE 正文或补充表格中体现。 \\
EC-34 & 风险成本 & 围绕“风险成本”说明其与建筑能源管理系统的关系，并给出成本、收益或管理价值判断。 & SEE 正文或补充表格中体现。 \\
EC-35 & 经济结论 & 围绕“经济结论”说明其与建筑能源管理系统的关系，并给出成本、收益或管理价值判断。 & SEE 正文或补充表格中体现。 \\
\bottomrule
\end{longtable}

\subsection{财务分析展开项}
财务分析展开项用于说明本项目的经济分析并非脱离系统泛泛而谈，而是与每个功能模块、成本来源和收益来源直接关联。
\begin{longtable}{L{1.7cm}L{4.0cm}L{6.4cm}L{3.4cm}}
\caption{财务分析展开项}\\
\toprule
编号 & 分析项 & 说明 & 经济意义 \\
\midrule
\endfirsthead
\toprule
编号 & 分析项 & 说明 & 经济意义 \\
\midrule
\endhead
FA-01 & 开发投入折算 & 对“开发投入折算”进行成本或收益解释，并说明它如何影响项目投入、维护费用、节能收益或管理价值。 & 用于支撑 SEE 结论。 \\
FA-02 & 测试投入折算 & 对“测试投入折算”进行成本或收益解释，并说明它如何影响项目投入、维护费用、节能收益或管理价值。 & 用于支撑 SEE 结论。 \\
FA-03 & 文档投入折算 & 对“文档投入折算”进行成本或收益解释，并说明它如何影响项目投入、维护费用、节能收益或管理价值。 & 用于支撑 SEE 结论。 \\
FA-04 & 维护投入折算 & 对“维护投入折算”进行成本或收益解释，并说明它如何影响项目投入、维护费用、节能收益或管理价值。 & 用于支撑 SEE 结论。 \\
FA-05 & 人员机会成本 & 对“人员机会成本”进行成本或收益解释，并说明它如何影响项目投入、维护费用、节能收益或管理价值。 & 用于支撑 SEE 结论。 \\
FA-06 & 本地部署节省 & 对“本地部署节省”进行成本或收益解释，并说明它如何影响项目投入、维护费用、节能收益或管理价值。 & 用于支撑 SEE 结论。 \\
FA-07 & 云部署备选成本 & 对“云部署备选成本”进行成本或收益解释，并说明它如何影响项目投入、维护费用、节能收益或管理价值。 & 用于支撑 SEE 结论。 \\
FA-08 & 模型调用边际成本 & 对“模型调用边际成本”进行成本或收益解释，并说明它如何影响项目投入、维护费用、节能收益或管理价值。 & 用于支撑 SEE 结论。 \\
FA-09 & 数据接入追加成本 & 对“数据接入追加成本”进行成本或收益解释，并说明它如何影响项目投入、维护费用、节能收益或管理价值。 & 用于支撑 SEE 结论。 \\
FA-10 & 真实设备集成成本 & 对“真实设备集成成本”进行成本或收益解释，并说明它如何影响项目投入、维护费用、节能收益或管理价值。 & 用于支撑 SEE 结论。 \\
FA-11 & 权限系统追加成本 & 对“权限系统追加成本”进行成本或收益解释，并说明它如何影响项目投入、维护费用、节能收益或管理价值。 & 用于支撑 SEE 结论。 \\
FA-12 & 数据库替换成本 & 对“数据库替换成本”进行成本或收益解释，并说明它如何影响项目投入、维护费用、节能收益或管理价值。 & 用于支撑 SEE 结论。 \\
FA-13 & 多校区推广成本 & 对“多校区推广成本”进行成本或收益解释，并说明它如何影响项目投入、维护费用、节能收益或管理价值。 & 用于支撑 SEE 结论。 \\
FA-14 & 运维培训成本 & 对“运维培训成本”进行成本或收益解释，并说明它如何影响项目投入、维护费用、节能收益或管理价值。 & 用于支撑 SEE 结论。 \\
FA-15 & 使用阻力成本 & 对“使用阻力成本”进行成本或收益解释，并说明它如何影响项目投入、维护费用、节能收益或管理价值。 & 用于支撑 SEE 结论。 \\
FA-16 & 低效设备发现收益 & 对“低效设备发现收益”进行成本或收益解释，并说明它如何影响项目投入、维护费用、节能收益或管理价值。 & 用于支撑 SEE 结论。 \\
FA-17 & 夜间负荷治理收益 & 对“夜间负荷治理收益”进行成本或收益解释，并说明它如何影响项目投入、维护费用、节能收益或管理价值。 & 用于支撑 SEE 结论。 \\
FA-18 & COP 优化收益 & 对“COP 优化收益”进行成本或收益解释，并说明它如何影响项目投入、维护费用、节能收益或管理价值。 & 用于支撑 SEE 结论。 \\
FA-19 & 人工报表节约 & 对“人工报表节约”进行成本或收益解释，并说明它如何影响项目投入、维护费用、节能收益或管理价值。 & 用于支撑 SEE 结论。 \\
FA-20 & 异常闭环管理价值 & 对“异常闭环管理价值”进行成本或收益解释，并说明它如何影响项目投入、维护费用、节能收益或管理价值。 & 用于支撑 SEE 结论。 \\
FA-21 & 数据口径统一价值 & 对“数据口径统一价值”进行成本或收益解释，并说明它如何影响项目投入、维护费用、节能收益或管理价值。 & 用于支撑 SEE 结论。 \\
FA-22 & 管理可视化价值 & 对“管理可视化价值”进行成本或收益解释，并说明它如何影响项目投入、维护费用、节能收益或管理价值。 & 用于支撑 SEE 结论。 \\
FA-23 & 系统展示价值 & 对“系统展示价值”进行成本或收益解释，并说明它如何影响项目投入、维护费用、节能收益或管理价值。 & 用于支撑 SEE 结论。 \\
FA-24 & 课程学习价值 & 对“课程学习价值”进行成本或收益解释，并说明它如何影响项目投入、维护费用、节能收益或管理价值。 & 用于支撑 SEE 结论。 \\
FA-25 & 试点推广价值 & 对“试点推广价值”进行成本或收益解释，并说明它如何影响项目投入、维护费用、节能收益或管理价值。 & 用于支撑 SEE 结论。 \\
FA-26 & 保守 ROI & 对“保守 ROI”进行成本或收益解释，并说明它如何影响项目投入、维护费用、节能收益或管理价值。 & 用于支撑 SEE 结论。 \\
FA-27 & 基准 ROI & 对“基准 ROI”进行成本或收益解释，并说明它如何影响项目投入、维护费用、节能收益或管理价值。 & 用于支撑 SEE 结论。 \\
FA-28 & 乐观 ROI & 对“乐观 ROI”进行成本或收益解释，并说明它如何影响项目投入、维护费用、节能收益或管理价值。 & 用于支撑 SEE 结论。 \\
FA-29 & 回收期 & 对“回收期”进行成本或收益解释，并说明它如何影响项目投入、维护费用、节能收益或管理价值。 & 用于支撑 SEE 结论。 \\
FA-30 & 盈亏平衡点 & 对“盈亏平衡点”进行成本或收益解释，并说明它如何影响项目投入、维护费用、节能收益或管理价值。 & 用于支撑 SEE 结论。 \\
\bottomrule
\end{longtable}

\end{document}
